\documentclass{beamer}

\usepackage{amsmath,amsfonts,amssymb,bm}
\usepackage{quantikz}
\usepackage{graphicx}

\usetheme{Madrid}

\title{Analog and Digital Quantum Computing}
\subtitle{From Hamiltonian Dynamics to Variational Algorithms}
\author{Morten Hjorth-Jensen}
\date{Spring 2026}

\begin{document}

\frame{\titlepage}

%================================================
\section{Motivation}
%================================================

\begin{frame}{Quantum Computing as Physics}
\[
i\hbar \frac{d}{dt} |\psi\rangle = H |\psi\rangle
\]
\[
|\psi(t)\rangle = e^{-iHt}|\psi(0)\rangle
\]

\begin{itemize}
\item Computation = control of quantum dynamics
\end{itemize}
\end{frame}

%------------------------------------------------

\begin{frame}{Two Paradigms}
\begin{itemize}
\item Digital quantum computing
\item Analog quantum computing
\end{itemize}

\begin{block}{Key distinction}
Discrete gate decomposition vs continuous evolution
\end{block}
\end{frame}

%================================================
\section{Digital Quantum Computing}
%================================================

\begin{frame}{Gate-Based Model}
\[
U = \prod_i U_i
\]
\begin{itemize}
\item Universal computation
\end{itemize}
\end{frame}

%------------------------------------------------

\begin{frame}{Quantum Circuit Example}
\[
\begin{quantikz}
\lstick{|0\rangle} & \gate{H} & \ctrl{1} & \meter{} \\
\lstick{|0\rangle} & \qw & \targ{} & \meter{}
\end{quantikz}
\]
\end{frame}

%------------------------------------------------

\begin{frame}{Trotterization}
\[
e^{-i(H_1+H_2)t} \approx
\left(e^{-iH_1\Delta t}e^{-iH_2\Delta t}\right)^n
\]
\end{frame}

%------------------------------------------------

\begin{frame}{Example: QAOA}
\[
|\psi\rangle =
\prod_k e^{-i\beta_k H_M} e^{-i\gamma_k H_C}|+\rangle
\]
\end{frame}

%------------------------------------------------

\begin{frame}{Example: HHL}
\begin{itemize}
\item Spectral decomposition
\item Implements $A^{-1}$
\end{itemize}
\end{frame}

%------------------------------------------------

\begin{frame}{Advantages}
\begin{itemize}
\item Universal
\item Programmable
\item Error correction possible
\end{itemize}
\end{frame}

%------------------------------------------------

\begin{frame}{Limitations}
\begin{itemize}
\item Deep circuits
\item Noise
\item Trotter errors
\end{itemize}
\end{frame}

%================================================
\section{Analog Quantum Computing}
%================================================

\begin{frame}{Analog Model}
\[
|\psi(t)\rangle = e^{-iH_{\text{sim}}t}|\psi(0)\rangle
\]
\end{frame}

%------------------------------------------------

\begin{frame}{Simulation Principle}
\[
H_{\text{sim}} \approx H_{\text{target}}
\]
\end{frame}

%------------------------------------------------

\begin{frame}{Physical Platforms}
\begin{itemize}
\item Cold atoms
\item Trapped ions
\item Rydberg atoms
\end{itemize}
\end{frame}

%------------------------------------------------

\begin{frame}{Quantum Annealing}
\[
H(t) = (1-s)H_M + sH_C
\]
\end{frame}

%------------------------------------------------

\begin{frame}{Advantages}
\begin{itemize}
\item Natural dynamics
\item Efficient for many-body systems
\end{itemize}
\end{frame}

%------------------------------------------------

\begin{frame}{Limitations}
\begin{itemize}
\item Not universal
\item Limited control
\end{itemize}
\end{frame}

%================================================
\section{Analog Simulation of Many-Body Models}
%================================================

\begin{frame}{Ising Model}
\[
H = -J \sum_{ij} Z_i Z_j - h \sum_i X_i
\]
\begin{itemize}
\item Implemented in trapped ions, Rydberg systems
\end{itemize}
\end{frame}

%------------------------------------------------

\begin{frame}{Hubbard Model}
\[
H = -t \sum_{\langle ij\rangle} c_i^\dagger c_j
+ U \sum_i n_{i\uparrow} n_{i\downarrow}
\]
\begin{itemize}
\item Realized in optical lattices
\end{itemize}
\end{frame}

%------------------------------------------------

\begin{frame}{Why Analog is Powerful}
\begin{itemize}
\item Direct access to many-body dynamics
\item No exponential classical cost
\end{itemize}
\end{frame}

%================================================
\section{Connections to Linear Response and TDHF}
%================================================

\begin{frame}{Linear Response Equation}
\[
(\omega I - M)x = b
\]
\end{frame}

%------------------------------------------------

\begin{frame}{TDHF Equation}
\[
i \frac{d\rho}{dt} = [h[\rho],\rho]
\]
\end{frame}

%------------------------------------------------

\begin{frame}{Linearized TDHF}
\[
(\omega I - \mathcal{L})\delta\rho = s
\]
\end{frame}

%------------------------------------------------

\begin{frame}{Connection to HHL}
\begin{itemize}
\item Same mathematical structure:
\[
A x = b
\]
\item Inverse operator = response function
\end{itemize}
\end{frame}

%================================================
\section{Variational Algorithms}
%================================================

\begin{frame}{Variational Principle}
\[
E = \langle \psi(\theta)|H|\psi(\theta)\rangle
\]
\end{frame}

%------------------------------------------------

\begin{frame}{ADAPT-VQE}
\begin{itemize}
\item Adaptive ansatz construction
\item Gradient selection:
\[
\langle [H,A_k] \rangle
\]
\end{itemize}
\end{frame}

%------------------------------------------------

\begin{frame}{Connection to Coupled Cluster}
\begin{itemize}
\item BCH expansion
\item Many-body correlations
\end{itemize}
\end{frame}

%------------------------------------------------

\begin{frame}{Variational Dynamics}
\begin{itemize}
\item McLachlan principle
\[
\delta \| (i\partial_t - H)|\psi\rangle \| = 0
\]
\end{itemize}
\end{frame}

%================================================
\section{Quantum Control Theory}
%================================================

\begin{frame}{Control Hamiltonian}
\[
H(t) = \sum_k u_k(t) H_k
\]
\end{frame}

%------------------------------------------------

\begin{frame}{Optimal Control}
\begin{itemize}
\item Find $u_k(t)$ to achieve target state
\end{itemize}
\end{frame}

%------------------------------------------------

\begin{frame}{Digital vs Analog Control}
\begin{itemize}
\item Digital: discrete pulses
\item Analog: continuous control
\end{itemize}
\end{frame}

%------------------------------------------------

\begin{frame}{QAOA as Control}
\begin{itemize}
\item Piecewise constant control protocol
\end{itemize}
\end{frame}

%================================================
\section{Unified View}
%================================================

\begin{frame}{Unified Framework}
\begin{itemize}
\item Analog: continuous evolution
\item Digital: discretized evolution
\item Variational: optimized evolution
\end{itemize}
\end{frame}

%------------------------------------------------

\begin{frame}{Core Mathematical Object}
\[
e^{-iHt}
\]
\begin{itemize}
\item Everything reduces to controlling this operator
\end{itemize}
\end{frame}

%------------------------------------------------

\begin{frame}{Conceptual Chain}
\begin{itemize}
\item Analog simulation → real dynamics
\item Digital circuits → approximated dynamics
\item Variational methods → optimized dynamics
\end{itemize}
\end{frame}

%================================================
\section{Outlook}
%================================================

\begin{frame}{Future Directions}
\begin{itemize}
\item Hybrid analog-digital systems
\item Quantum simulation of materials
\item Quantum-enhanced many-body theory
\end{itemize}
\end{frame}

%------------------------------------------------

\begin{frame}{Summary}
\begin{itemize}
\item Digital: universal and programmable
\item Analog: efficient and physics-driven
\item Variational: bridge between both
\end{itemize}
\end{frame}

\end{document}
